%% LyX 2.4.4 created this file.  For more info, see https://www.lyx.org/.
%% Do not edit unless you really know what you are doing.
\documentclass[english,t]{beamer}
\usepackage{mathptmx}
\usepackage[T1]{fontenc}
\usepackage[latin9]{inputenc}
\usepackage{amssymb}
\usepackage{cancel}
\PassOptionsToPackage{normalem}{ulem}
\usepackage{ulem}

\makeatletter

%%%%%%%%%%%%%%%%%%%%%%%%%%%%%% LyX specific LaTeX commands.
%% Because html converters don't know tabularnewline
\providecommand{\tabularnewline}{\\}

%%%%%%%%%%%%%%%%%%%%%%%%%%%%%% Textclass specific LaTeX commands.
% this default might be overridden by plain title style
\newcommand\makebeamertitle{\frame{\maketitle}}%
% (ERT) argument for the TOC
\AtBeginDocument{%
  \let\origtableofcontents=\tableofcontents
  \def\tableofcontents{\@ifnextchar[{\origtableofcontents}{\gobbletableofcontents}}
  \def\gobbletableofcontents#1{\origtableofcontents}
}

%%%%%%%%%%%%%%%%%%%%%%%%%%%%%% User specified LaTeX commands.
\usetheme{Warsaw}
% or ...

\setbeamercovered{transparent}
% or whatever (possibly just delete it)
\setbeamercovered{invisible} %makes overlays invisible


%gets rid of navigation symbols
\setbeamertemplate{navigation symbols}{}

\usepackage{multicol}

\usepackage{tikz}
\usetikzlibrary{calc}
\usetikzlibrary{plotmarks}
\usepackage{pgfplots}
\pgfplotsset{compat=newest}

\usepackage{calc}
\usepackage[nomessages]{fp}

\usepackage{advdate}    % Advancing/saving dates
\usepackage{datetime}   % Dates formatting
\usepackage{datenumber} % Counters for dates
\usepackage[super]{nth}

%\usepackage{pgfpages}
%\pgfpagesuselayout{4 on 1}[a4paper, landscape, border shrink=7mm]
%\pgfpageslogicalpageoptions{1}{border code=\pgfusepath{stroke}}
%\pgfpageslogicalpageoptions{2}{border code=\pgfusepath{stroke}}
%\pgfpageslogicalpageoptions{3}{border code=\pgfusepath{stroke}}
%\pgfpageslogicalpageoptions{4}{border code=\pgfusepath{stroke}}
%%%Don't forget to add 'handout'

\makeatother

\usepackage{babel}
\begin{document}
\newcommand{\xmax}{2000}
\newcommand{\xmin}{0}
\newcommand{\xstep}{0} %must be xmin+n
\newcommand{\ymax}{500}
\newcommand{\ymin}{0}
\newcommand{\ystep}{0} %must be ymin+n

\newcounter{countEx} \setcounter{countEx}{0}
\newcounter{countProb} \setcounter{countProb}{0}
\resetcounteronoverlays{countEx} \resetcounteronoverlays{countProb}

\newcommand{\ex}[3]{
	\addtocounter{countEx}{1}
	\only<#1-#2>{\begin{exampleblock} {Example \chNum.\secNum.\arabic{countEx}.} #3	\end{exampleblock}}
}

\newcommand{\prob}[4]{
	\addtocounter{countProb}{1}
	\only<#1-#2>{\begin{exampleblock} {Problem  \chNum.\secNum.\arabic{countProb}.\,#3} #4	\end{exampleblock}}
}

\newcommand{\yline}[4]{
	(#2 - #4)/(#1 - #3)*(x - #1) + #2
}

\beamerdefaultoverlayspecification{<+->}

%%%%Chapter and Section numbers%%%%%
\newcommand{\chNum}{6}
\newcommand{\secNum}{1}
\newcommand{\secTitle}{Statements}

\title[Section \chNum.\secNum: \secTitle]{Section \chNum.\secNum: \secTitle}

\subtitle{Finite Mathematics~~MTH 132~~Spring 2014}

\author{Dr.\,James Ashe}

\titlegraphic{\includegraphics[height=1.5cm]{/home/jim/JamesAshe/JCSU/images/jcsu_bull}}

\institute{Johnson C. Smith University}

\date{{}}

\makebeamertitle 
\begin{frame}{Chapter 6: Logic}

\begin{columns}


\column{11cm}
\begin{itemize}
\item []Chapter $6$: Logic

\begin{itemize}
\item [6.1:] Statements
\item [6.2:] Truth Tables and Equivalent Statements
\item [6.4:] Conditional Statements
\end{itemize}
\end{itemize}
\end{columns}
\end{frame}

\begin{frame}{}

\begin{block}{Statement}
A \emph{statement }is a declarative sentence that is either true
or false. 
\end{block}
\begin{columns}


\column{5.5cm}
\begin{itemize}
\item <2->Socrates was a man. \uncover<6->{\textbf{T}}
\item <3->$8-3+4=1$ \uncover<7->{\textbf{F}}
\item <4->A triangle has three sides. \uncover<8->{\textbf{T}}
\item <5->[]\uline{STATEMENTS: True or False}
\end{itemize}

\column{5.5cm}
\begin{itemize}
\item <9->Is Socrates a man?
\item <10->Go do some math.
\item <11->$8-3+4$
\item <12->This statement is false (\emph{Liar's Paradox}). 
\item <13->$A$ and not $A$.
\item <14->[]\uline{NOT STATEMENTS}
\end{itemize}
\end{columns}
\end{frame}

\begin{frame}{}

\begin{block}{Compound Statement}
A \emph{compound statement} is two or more statements combined.

\begin{itemize}
\item <2->Combined with \emph{logical connectives}, e.g.,\emph{ }and, or,
not, yet, if, then,...
\item <3->The simple statements that compose a compound statement are called
\emph{component} statements.
\end{itemize}
\end{block}
\begin{columns}


\column{5.5cm}
\begin{itemize}
\item <4->Socrates was a man, and Theano was a woman. 
\item <5->If a polygon has three sides then it is a triangle.
\item <6->$8-3+4=9$ not $1$.
\item <7->[]\uline{COMPOUND STATEMENTS}
\end{itemize}

\column{5.5cm}
\begin{itemize}
\item <8->Is Socrates a man or woman?
\item <9->If a polygon has three sides than go get some ice cream.
\item <10->Truth or dare is a game.
\item <11->This statement is true, and this statement is false. 
\item <12->[]\uline{NOT COMPOUND STATEMENTS}
\end{itemize}
\end{columns}
\end{frame}

\begin{frame}{Negation}

\emph{``No agreement exists as to the possibility of defining negation,
as to its logical status, function, and meaning, as to its field of
applicability..., and as to the interpretation of the negative judgment''
}(F.H. Heinemann 1944)
\begin{block}<2->{Negation}
The \emph{negation }of a true statement makes it false, and the
\emph{negation }of a false statement is true.
\end{block}
\begin{columns}


\column{5.5cm}
\begin{itemize}
\item <2->[]\textbf{Statement}

\begin{itemize}
\item <2->Socrates is a man. 
\item <4->Socrates is not a man.
\item <6->$3<4$, ``$3$ is less than $4$'' 
\item <8->$8-3+4=9$
\end{itemize}
\end{itemize}

\column{6.5cm}
\begin{itemize}
\item <3->[]\textbf{Negation of Statement}

\begin{itemize}
\item <3->Socrates is not a man?
\item <5->Socrates is a man.
\item <7->$3\geq4$, ``$3$ is not less than $4$''
\item <9->$8-3+4\neq9$
\end{itemize}
\end{itemize}
\end{columns}
\end{frame}

\begin{frame}{Symbolic Logic}

\begin{columns}


\column{4.25cm}
\begin{block}{Notation}

\begin{tabular}{rl}
 & \textbf{Symbol}\tabularnewline
statement & $p$, $q$, $r$, etc.\tabularnewline
``and'' & $\wedge$\tabularnewline
``or'' & $\vee$\tabularnewline
``not'' & $\sim$\tabularnewline
\end{tabular}
\end{block}

\column{6.5cm}

\uncover<2->{Let $m$ represent ``Socrates is a man'' and $r$
represent ``it's going to rain''.}
\begin{itemize}
\item <2->[ex.]$m\wedge r$

\begin{itemize}
\item []<3->\emph{Socrates is a man and it's going to rain.}
\end{itemize}
\item <4->[ex.]$\sim m\vee r$

\begin{itemize}
\item []<5->\emph{Socrates if }not \emph{a man or it's going to rain.}
\end{itemize}
\item <6->[ex.]$\sim(m\vee r)$

\begin{itemize}
\item []<7->\emph{It is not case that Socrates is a man or it's going to
rain.}
\item []<8->Equivalently, $\sim m\wedge\sim r$, \emph{Socrates is not
a man and it is not going to rain}.
\end{itemize}
\end{itemize}
\end{columns}
\end{frame}

\begin{frame}{Your Turn: Statements}

\prob{1}{1}{Decide whether the following is a statement. If it is a statement, decide whether or not it is compound}{Nashville is the capital of Tennessee.}

\prob{2}{2}{Decide whether the following is a statement. If it is a statement, decide whether or not it is compound}{Got milk?}

\prob{3}{3}{Decide whether the following is a statement. If it is a statement, decide whether or not it is compound}{China does not have a population of more than 1 billion, and it is a country.}

\prob{4}{4}{Decide whether the following is a statement. If it is a statement, decide whether or not it is compound}{If I get an \emph{A}  then I will celebrate.}
\end{frame}

\begin{frame}{Your Turn: Negation}

\prob{1}{1}{Write a negation of each statement.}{This is not the time to complain.}\prob{2}{2}{Write a negation of each statement.}{$y\leq 5$}\prob{3}{3}{Write a negation of each statement.}{$p \wedge q$}
\end{frame}

\begin{frame}{Finding the Truth Value: Conjunction, i.e., ``and''}

\ex{1}{}{Let $p$ represent ``$5>3$'' and $q$ represent ``$6<0$''. Find the truth value of $p \wedge q$}
\begin{enumerate}
\item <2->[]Using a \emph{truth table}. 

\begin{enumerate}
\item <3->[]%
\begin{tabular}{cc|c}
$p$ & $q$ & $p\wedge q$\tabularnewline
\hline 
 &  & \tabularnewline
 &  & \tabularnewline
 &  & \tabularnewline
 &  & \tabularnewline
\end{tabular}
\end{enumerate}
\begin{itemize}
\item <4->Determine which statements are true and false, and choose the
corresponding answer.
\end{itemize}
\end{enumerate}
\end{frame}

\begin{frame}{Finding the Truth Value: Conjunction, i.e., ``and''}

\ex{1}{}{Let $p$ be true and $q$ be true. Find the truth value of $p \wedge q$}
\begin{enumerate}
\item <2->[]Using a \emph{truth table}. 

\begin{enumerate}
\item <3->[]%
\begin{tabular}{cc|c}
$p$ & $q$ & $p\wedge q$\tabularnewline
\hline 
T & T & T\tabularnewline
T & F & F\tabularnewline
F & T & F\tabularnewline
F & F & F\tabularnewline
\end{tabular}
\end{enumerate}
\begin{itemize}
\item <4->Determine which statements are true and false, and choose the
corresponding answer.
\end{itemize}
\end{enumerate}
\end{frame}

\begin{frame}{Finding the Truth Value: Disjunction, i.e., ``or''}

\ex{1}{}{Let $p$ represent ``$5>3$'' or $q$ represent ``$6<0$''. Find the truth value of $p \wedge q$}
\begin{enumerate}
\item <2->[]Using a \emph{truth table}. 

\begin{enumerate}
\item <3->[]%
\begin{tabular}{cc|c}
$p$ & $q$ & $p\vee q$\tabularnewline
\hline 
T & T & \tabularnewline
T & F & \tabularnewline
F & T & \tabularnewline
F & F & \tabularnewline
\end{tabular}
\end{enumerate}
\begin{itemize}
\item <4->Determine which statements are true and false, and choose the
corresponding answer.
\end{itemize}
\end{enumerate}
\end{frame}

\begin{frame}{Compound Statements}

\ex{1}{}{Let $p$ represent a false statement, and let $q$ represent a true statement. Find the truth value of the compound statement $\sim (\sim p \wedge q)$}
\begin{columns}


\column{4.5cm}
\begin{enumerate}
\item <2->[]\footnotesize%
\begin{tabular}{cc|c}
$p$ & $q$ & $p\wedge q$\tabularnewline
\hline 
T & T & T\tabularnewline
T & F & F\tabularnewline
F & T & F\tabularnewline
F & F & F\tabularnewline
\end{tabular}
\item <3->[]\footnotesize%
\begin{tabular}{cc|c}
 & \multicolumn{1}{c}{} & \tabularnewline
$p$ & $q$ & $p\vee q$\tabularnewline
\hline 
T & T & T\tabularnewline
T & F & T\tabularnewline
F & T & T\tabularnewline
F & F & F\tabularnewline
\end{tabular}
\end{enumerate}

\column{6.5cm}
\begin{itemize}
\item []<4->\large$\sim(\sim p\wedge q)$
\end{itemize}
\end{columns}
\end{frame}

\begin{frame}{Compound Statements}

\ex{1}{}{Let $p$ represent a true statement, and let $q$ and $r$ represent false statements. Find the truth value of the compound statement $(p \wedge r) \vee \sim q$}
\begin{columns}


\column{4.5cm}
\begin{enumerate}
\item <2->[]\footnotesize%
\begin{tabular}{cc|c}
$p$ & $q$ & $p\wedge q$\tabularnewline
\hline 
T & T & T\tabularnewline
T & F & F\tabularnewline
F & T & F\tabularnewline
F & F & F\tabularnewline
\end{tabular}
\item <3->[]\footnotesize%
\begin{tabular}{cc|c}
 & \multicolumn{1}{c}{} & \tabularnewline
$p$ & $q$ & $p\vee q$\tabularnewline
\hline 
T & T & T\tabularnewline
T & F & T\tabularnewline
F & T & T\tabularnewline
F & F & F\tabularnewline
\end{tabular}
\end{enumerate}

\column{6.5cm}
\begin{itemize}
\item []<4->\large$(p\wedge r)\vee\sim q$
\end{itemize}
\end{columns}
\end{frame}

\begin{frame}{Finding the Truth Value: Compound Statements}

\ex{1}{}{If $q$ is true, what is the truth value of $q \vee (q \wedge \sim p)$}
\begin{columns}


\column{4.5cm}
\begin{enumerate}
\item <2->[]\footnotesize%
\begin{tabular}{cc|c}
$p$ & $q$ & $p\wedge q$\tabularnewline
\hline 
T & T & T\tabularnewline
T & F & F\tabularnewline
F & T & F\tabularnewline
F & F & F\tabularnewline
\end{tabular}
\item <3->[]\footnotesize%
\begin{tabular}{cc|c}
 & \multicolumn{1}{c}{} & \tabularnewline
$p$ & $q$ & $p\vee q$\tabularnewline
\hline 
T & T & T\tabularnewline
T & F & T\tabularnewline
F & T & T\tabularnewline
F & F & F\tabularnewline
\end{tabular}
\end{enumerate}

\column{6.5cm}
\begin{itemize}
\item []<4->\large$q\vee(q\wedge\sim p)$
\end{itemize}
\end{columns}
\end{frame}

\begin{frame}{Your Turn: Find the Truth Value}

\prob{1}{1}{Let $p$ represent a false statement and $q$ represent a true statement. Find the truth value of each compound statement.}{$\sim p$} 

\prob{2}{2}{Let $p$ represent a false statement and $q$ represent a true statement. Find the truth value of each compound statement.}{$p \vee \sim q$} 

\prob{3}{3}{Let $p$ represent a false statement and $q$ represent a true statement. Find the truth value of each compound statement.}{$\sim (p \vee \sim q)$} 
\end{frame}

\begin{frame}{Your Turn: Find the Truth Value}

\prob{1}{1}{Let $p$ represent a true statement, and $q$ and $r$  represent false statements. Find the truth value of each compound statement.}{$(q \vee \sim r) \wedge p$} 

\prob{2}{2}{Let $p$ represent a true statement, and $q$ and $r$  represent false statements. Find the truth value of each compound statement.}{$(\sim p \wedge q) \vee \sim r$} 
\end{frame}

\begin{frame}{Your Turn: Find the Truth Value}

\prob{1}{1}{Let $p$ be the statement``$2>7$'', $q$ be statement ``$8 \leq 6$,'' and $r$ be the statement ``$19 \leq 19$.'' Find the truth value of each compound statement.}{$\sim p \vee \sim r$}

\prob{2}{2}{Let $p$ be the statement``$2>7$'', $q$ be statement ``$8 \leq 6$,'' and $r$ be the statement ``$19 \leq 19$.'' Find the truth value of each compound statement.}{$(p \wedge q) \vee r$}

\prob{3}{3}{Let $p$ be the statement``$2>7$'', $q$ be statement ``$8 \leq 6$,'' and $r$ be the statement ``$19 \leq 19$.'' Find the truth value of each compound statement.}{$(\sim r \wedge q) \vee \sim p$}
\end{frame}

\begin{frame}{Answers}

Answers:\textbf{ T, F, T, T, T, F, T, T}
\end{frame}

\end{document}
